\renewcommand{\thesection}{\arabic{section}}\renewcommand{\theHsection}{chapter.\arabic{section}}\setcounter{section}{3}\section{不定积分}
\noindent\begin{minipage}{\textwidth}\ProblemHeading{例题}{4.1}
求不定积分 $\displaystyle\int \frac{x^2-3x+1}{x+1}\,dx$。
\end{minipage}\par\vfill
\noindent\begin{minipage}{\textwidth}\ProblemHeading{例题}{4.2}
求不定积分 $\displaystyle\int \frac1{\sin^2x\cos^2x}\,dx$。
\end{minipage}\par\vfill
\noindent\begin{minipage}{\textwidth}\ProblemHeading{例题}{4.3}
求不定积分 $\displaystyle\int \frac1{x^2(1+x^2)}\,dx$。
\end{minipage}\par\vfill
\noindent\begin{minipage}{\textwidth}\ProblemHeading{例题}{4.4}
求不定积分 $\displaystyle\int \tan x\,dx$。
\end{minipage}\par\vfill
\newpage
\noindent\begin{minipage}{\textwidth}\ProblemHeading{例题}{4.5}
求不定积分 $\displaystyle\int \frac{\ln x}{x\sqrt{1+\ln x}}\,dx$。
\end{minipage}\par\vfill
\noindent\begin{minipage}{\textwidth}\ProblemHeading{例题}{4.6}
求不定积分 $\displaystyle\int \frac1{\sqrt{e^x+1}}\,dx$。
\end{minipage}\par\vfill
\noindent\begin{minipage}{\textwidth}\ProblemHeading{例题}{4.7}
求不定积分 $\displaystyle\int \frac1{\sqrt{x}(1+\sqrt[3]{x})}\,dx$。
\end{minipage}\par\vfill
\noindent\begin{minipage}{\textwidth}\ProblemHeading{例题}{4.8}
求不定积分 $\displaystyle\int \sqrt{a^2-x^2}\,dx$。其中 $a>0$ 为常数。
\end{minipage}\par\vfill
\newpage
\noindent\begin{minipage}{\textwidth}\ProblemHeading{例题}{4.9}
求不定积分 $\displaystyle\int \frac1{\sqrt{x^2-a^2}}\,dx$。其中 $a>0$ 为常数。
\end{minipage}\par\vfill
\noindent\begin{minipage}{\textwidth}\ProblemHeading{例题}{4.10}
求不定积分 $\displaystyle\int \frac1{x\sqrt{a^2-x^2}}\,dx$。其中 $a>0$ 为常数。
\end{minipage}\par\vfill
\noindent\begin{minipage}{\textwidth}\ProblemHeading{例题}{4.11}
求 $\displaystyle\int\frac{x^2+1}{x^4+x^2+1}\,dx$，要求所得原函数在整个实轴上成立。
\end{minipage}\par\vfill
\noindent\begin{minipage}{\textwidth}\ProblemHeading{例题}{4.12}
求不定积分 $\displaystyle\int \ln x\,dx$。
\end{minipage}\par\vfill
\newpage
\noindent\begin{minipage}{\textwidth}\ProblemHeading{例题}{4.13}
求不定积分 $\displaystyle\int e^{ax}\sin bx\,dx$。其中 $a,b$ 为非零实数。
\end{minipage}\par\vfill
\noindent\begin{minipage}{\textwidth}\ProblemHeading{例题}{4.14}
记 $I_n=\displaystyle\int\frac1{(x^2+a^2)^n}\,dx$，其中 $a\ne0$ 为实数，$n=1,2,\ldots$。证明递推公式\[I_{n+1}=\frac{x}{2na^2(x^2+a^2)^n}+\frac{2n-1}{2na^2}I_n.\]
\end{minipage}\par\vfill
\noindent\begin{minipage}{\textwidth}\ProblemHeading{例题}{4.15}
求不定积分 $\displaystyle\int xe^x(1+e^x)^{-3/2}\,dx$。
\end{minipage}\par\vfill
\noindent\begin{minipage}{\textwidth}\ProblemHeading{习题}{4.16}
求下列不定积分。
\SubQ{1}{$\displaystyle\int x(x-1)^3\,dx$。}
\SubQ{2}{$\displaystyle\int \frac{e^{3x}+1}{e^x+1}\,dx$。}
\SubQ{3}{$\displaystyle\int (2^x+3^x)^2\,dx$。}
\end{minipage}\par\vfill
\newpage
\noindent\begin{minipage}{\textwidth}\ProblemHeading{习题}{4.17}
求下列不定积分。
\SubQ{1}{$\displaystyle\int \tan^2x\,dx$。}
\SubQ{2}{$\displaystyle\int \frac{x^2}{1+x^2}\,dx$。}
\SubQ{3}{$\displaystyle\int \frac{1+\cos^2x}{1+\cos2x}\,dx$。}
\end{minipage}\par\vfill
\noindent\begin{minipage}{\textwidth}\ProblemHeading{习题}{4.18}
用第一代换法求下列不定积分。
\SubQ{1}{$\displaystyle\int (2x-1)^{100}\,dx$。}
\SubQ{2}{$\displaystyle\int \frac1{x^2}\sin\frac1x\,dx$。}
\SubQ{3}{$\displaystyle\int \frac{\cos x-\sin x}{1+\sin x+\cos x}\,dx$。}
\end{minipage}\par\vfill
\noindent\begin{minipage}{\textwidth}\ProblemHeading{习题}{4.19}
用第一代换法求下列不定积分。
\SubQ{1}{$\displaystyle\int \frac{\arctan x}{1+x^2}\,dx$。}
\SubQ{2}{$\displaystyle\int x\sqrt{1-x^2}\,dx$。}
\SubQ{3}{$\displaystyle\int \frac1{\sqrt{x}(1+x)}\,dx$。}
\end{minipage}\par\vfill
\noindent\begin{minipage}{\textwidth}\ProblemHeading{习题}{4.20}
用第一代换法求下列不定积分。
\SubQ{1}{$\displaystyle\int \frac{\arctan(1/x)}{1+x^2}\,dx$。}
\SubQ{2}{$\displaystyle\int \frac{1+\ln x}{1+x\ln x}\,dx$。}
\SubQ{3}{$\displaystyle\int \sin^2x\,dx$。}
\end{minipage}\par\vfill
\newpage
\noindent\begin{minipage}{\textwidth}\ProblemHeading{习题}{4.21}
用第一代换法求下列不定积分。
\SubQ{1}{$\displaystyle\int \sin^5x\cos x\,dx$。}
\end{minipage}\par\vfill
\noindent\begin{minipage}{\textwidth}\ProblemHeading{习题}{4.22}
用第二代换法求下列不定积分，其中 $a>0$ 为常数。
\SubQ{1}{$\displaystyle\int \sqrt{e^x-2}\,dx$。}
\SubQ{2}{$\displaystyle\int \sqrt{x^2+a^2}\,dx$。}
\SubQ{3}{$\displaystyle\int \frac1{(x^2-a^2)^{3/2}}\,dx$。}
\end{minipage}\par\vfill
\noindent\begin{minipage}{\textwidth}\ProblemHeading{习题}{4.23}
用第二代换法求下列不定积分，其中 $a>0$ 为常数。
\SubQ{1}{$\displaystyle\int \frac{x^2}{\sqrt{a^2-x^2}}\,dx$。}
\SubQ{2}{$\displaystyle\int \frac1{1+\sqrt{x+1}}\,dx$。}
\SubQ{3}{$\displaystyle\int \frac{x\ln x}{(1+x^2)^{3/2}}\,dx$。}
\end{minipage}\par\vfill
\noindent\begin{minipage}{\textwidth}\ProblemHeading{习题}{4.24}
用第二代换法求下列不定积分，其中 $a>0$ 为常数。
\SubQ{1}{$\displaystyle\int \frac{1-\ln x}{(x-\ln x)^2}\,dx$。}
\SubQ{2}{$\displaystyle\int \frac1{x^2\sqrt{x^2+a^2}}\,dx$。}
\SubQ{3}{$\displaystyle\int \frac{x+2}{\sqrt[3]{2x+1}}\,dx$。}
\end{minipage}\par\vfill
\newpage
\noindent\begin{minipage}{\textwidth}\ProblemHeading{习题}{4.25}
用第二代换法求下列不定积分，其中 $a>0$ 为常数。
\SubQ{1}{$\displaystyle\int \frac{x^{1/7}+x^{1/2}}{x^{8/7}+x^{1/14}}\,dx$。}
\SubQ{2}{$\displaystyle\int \frac{x-1}{x^2\sqrt{x^2-1}}\,dx$。}
\SubQ{3}{$\displaystyle\int \frac1{x^8(1+x^2)}\,dx$。}
\end{minipage}\par\vfill
\noindent\begin{minipage}{\textwidth}\ProblemHeading{习题}{4.26}
求下列不定积分。
\SubQ{1}{$\displaystyle\int |x|\,dx$。}
\SubQ{2}{$\displaystyle\int \max\{1,x^2\}\,dx$。}
\end{minipage}\par\vfill
\noindent\begin{minipage}{\textwidth}\ProblemHeading{习题}{4.27}
用分部积分法求下列不定积分。
\SubQ{1}{$\displaystyle\int x\sin x\,dx$。}
\SubQ{2}{$\displaystyle\int x^2\ln x\,dx$。}
\SubQ{3}{$\displaystyle\int \cos(\ln x)\,dx$。}
\end{minipage}\par\vfill
\noindent\begin{minipage}{\textwidth}\ProblemHeading{习题}{4.28}
用分部积分法求下列不定积分。
\SubQ{1}{$\displaystyle\int x^2\cos5x\,dx$。}
\SubQ{2}{$\displaystyle\int \sec^3x\,dx$。}
\SubQ{3}{$\displaystyle\int x^2e^x\,dx$。}
\end{minipage}\par\vfill
\newpage
\noindent\begin{minipage}{\textwidth}\ProblemHeading{习题}{4.29}
用分部积分法求下列不定积分。
\SubQ{1}{$\displaystyle\int x\arcsin x\,dx$。}
\SubQ{2}{$\displaystyle\int x(\arctan x)^2\,dx$。}
\SubQ{3}{$\displaystyle\int (\arcsin x)^2\,dx$。}
\end{minipage}\par\vfill
\noindent\begin{minipage}{\textwidth}\ProblemHeading{习题}{4.30}
用分部积分法求下列不定积分。
\SubQ{1}{$\displaystyle\int \ln(x+\sqrt{x^2+1})\,dx$。}
\end{minipage}\par\vfill
\noindent\begin{minipage}{\textwidth}\ProblemHeading{习题}{4.31}
导出下列不定积分的递推公式，其中 $n=1,2,\ldots$。
\SubQ{1}{$\displaystyle\int \sin^nx\,dx$。}
\SubQ{2}{$\displaystyle\int x^ne^x\,dx$。}
\end{minipage}\par\vfill
\noindent\begin{minipage}{\textwidth}\ProblemHeading{习题}{4.32}
求下列不定积分。
\SubQ{1}{$\displaystyle\int \frac1{1+e^x}\,dx$。}
\SubQ{2}{$\displaystyle\int \frac{x^2-1}{x^4+x^2+1}\,dx$。}
\SubQ{3}{$\displaystyle\int \frac1{x^4+x^6}\,dx$。}
\end{minipage}\par\vfill
\newpage
\noindent\begin{minipage}{\textwidth}\ProblemHeading{习题}{4.33}
求下列不定积分。
\SubQ{1}{$\displaystyle\int x\sqrt{x-2}\,dx$。}
\SubQ{2}{$\displaystyle\int \frac{\sqrt{x-1}\arctan\sqrt{x-1}}x\,dx$。}
\SubQ{3}{$\displaystyle\int \frac{xe^x}{\sqrt{e^x-2}}\,dx$。}
\end{minipage}\par\vfill
\noindent\begin{minipage}{\textwidth}\ProblemHeading{习题}{4.34}
求下列不定积分。
\SubQ{1}{$\displaystyle\int xe^x\sin x\,dx$。}
\SubQ{2}{$\displaystyle\int \frac1{(1+\tan x)\sin^2x}\,dx$。}
\SubQ{3}{$\displaystyle\int \frac{\sqrt{1-x}}{1-\sqrt{x}}\,dx$。}
\end{minipage}\par\vfill
\noindent\begin{minipage}{\textwidth}\ProblemHeading{习题}{4.35}
求下列不定积分。
\SubQ{1}{$\displaystyle\int \sqrt{\frac{x-1}{x+1}}\frac1{x^2}\,dx$。}
\SubQ{2}{$\displaystyle\int \frac{x\arctan x}{(1+x^2)^3}\,dx$。}
\SubQ{3}{$\displaystyle\int \frac{x}{1+\sin x}\,dx$。}
\end{minipage}\par\vfill
\noindent\begin{minipage}{\textwidth}\ProblemHeading{习题}{4.36}
求下列不定积分。
\SubQ{1}{$\displaystyle\int \arcsin\sqrt{x}\,dx$。}
\SubQ{2}{$\displaystyle\int \frac{x+\sin x}{1+\cos x}\,dx$。}
\SubQ{3}{$\displaystyle\int x\sin^2x\,dx$。}
\end{minipage}\par\vfill
\newpage
\noindent\begin{minipage}{\textwidth}\ProblemHeading{习题}{4.37}
求下列不定积分。
\SubQ{1}{$\displaystyle\int \frac{x^3}{\sqrt{1+x^2}}\,dx$。}
\SubQ{2}{$\displaystyle\int \frac{\arctan x}{x^2(1+x^2)}\,dx$。}
\SubQ{3}{$\displaystyle\int \frac{\arctan e^x}{e^x}\,dx$。}
\end{minipage}\par\vfill
\noindent\begin{minipage}{\textwidth}\ProblemHeading{习题}{4.38}
求下列不定积分。
\SubQ{1}{$\displaystyle\int e^{2x}(1+\tan x)^2\,dx$。}
\SubQ{2}{$\displaystyle\int \frac{x^2}{(x\sin x+\cos x)^2}\,dx$。}
\SubQ{3}{$\displaystyle\int \cos x\cos2x\cos3x\,dx$。}
\end{minipage}\par\vfill
\noindent\begin{minipage}{\textwidth}\ProblemHeading{习题}{4.39}
求下列不定积分。
\SubQ{1}{$\displaystyle\int \frac1{\sqrt{x-1}+\sqrt{x+1}}\,dx$。}
\SubQ{2}{$\displaystyle\int \frac1{\sqrt{\sqrt{x}+1}}\,dx$。}
\SubQ{3}{$\displaystyle\int \sqrt{\frac{x}{1-x\sqrt{x}}}\,dx$。}
\end{minipage}\par\vfill
\noindent\begin{minipage}{\textwidth}\ProblemHeading{习题}{4.40}
求下列不定积分。
\SubQ{1}{$\displaystyle\int e^{-x^2/2}\frac{\cos x-2x\sin x}{2\sqrt{\sin x}}\,dx$。}
\SubQ{2}{$\displaystyle\int \frac{xe^x}{(1+x)^2}\,dx$。}
\end{minipage}\par\vfill
\newpage
\noindent\begin{minipage}{\textwidth}\ProblemHeading{例题}{4.41}
求不定积分 $\displaystyle\int \frac1{x^3+1}\,dx$。
\end{minipage}\par\vfill
\noindent\begin{minipage}{\textwidth}\ProblemHeading{例题}{4.42}
求不定积分 $\displaystyle\int \frac1{x(x+1)(x^2+x+1)}\,dx$。
\end{minipage}\par\vfill
\noindent\begin{minipage}{\textwidth}\ProblemHeading{例题}{4.43}
求不定积分 $\displaystyle\int \frac1{x(x^8+1)}\,dx$。
\end{minipage}\par\vfill
\noindent\begin{minipage}{\textwidth}\ProblemHeading{例题}{4.44}
求不定积分 $\displaystyle\int \frac1{5+4\sin x}\,dx$。
\end{minipage}\par\vfill
\newpage
\noindent\begin{minipage}{\textwidth}\ProblemHeading{习题}{4.45}
设 $R(s,t)$ 为二元有理函数。利用代换 $u=\tanh(x/2)$，将 $\displaystyle\int R(\cosh x,\sinh x)\,dx$ 化为有理函数的不定积分。这里 $\cosh x=(e^x+e^{-x})/2$，$\sinh x=(e^x-e^{-x})/2$，$\tanh x=\sinh x/\cosh x$。
\end{minipage}\par\vfill
\noindent\begin{minipage}{\textwidth}\ProblemHeading{例题}{4.46}
求不定积分 $\displaystyle\int \sin^4x\,dx$。
\end{minipage}\par\vfill
\noindent\begin{minipage}{\textwidth}\ProblemHeading{习题}{4.47}
求下列有理函数的不定积分。
\SubQ{1}{$\displaystyle\int \frac1{x^2+x-2}\,dx$。}
\SubQ{2}{$\displaystyle\int \frac{x^4}{x^2+1}\,dx$。}
\SubQ{3}{$\displaystyle\int \frac{x^3+1}{x^3-x}\,dx$。}
\end{minipage}\par\vfill
\noindent\begin{minipage}{\textwidth}\ProblemHeading{习题}{4.48}
求下列有理函数的不定积分。
\SubQ{1}{$\displaystyle\int \frac1{(x^2+1)(x^2+x)}\,dx$。}
\SubQ{2}{$\displaystyle\int \frac{x}{(x+1)^2(x^2+x+1)}\,dx$。}
\SubQ{3}{$\displaystyle\int \frac{x^2+1}{x^4+1}\,dx$。}
\end{minipage}\par\vfill
\newpage
\noindent\begin{minipage}{\textwidth}\ProblemHeading{习题}{4.49}
求下列有理函数的不定积分。
\SubQ{1}{$\displaystyle\int \frac{x^5-x}{x^8+1}\,dx$。}
\SubQ{2}{$\displaystyle\int \frac{x^{15}}{(x^8+1)^2}\,dx$。}
\end{minipage}\par\vfill
\noindent\begin{minipage}{\textwidth}\ProblemHeading{习题}{4.50}
求下列三角函数有理式的不定积分，其中 $a,b$ 为常数，$a^2+b^2\ne0$。
\SubQ{1}{$\displaystyle\int \frac{1+\sin x}{\sin x(1+\cos x)}\,dx$。}
\SubQ{2}{$\displaystyle\int \frac{\sin^5x}{\cos x}\,dx$。}
\SubQ{3}{$\displaystyle\int \frac1{\sin^4x\cos^2x}\,dx$。}
\end{minipage}\par\vfill
\noindent\begin{minipage}{\textwidth}\ProblemHeading{习题}{4.51}
求下列三角函数有理式的不定积分，其中 $a,b$ 为常数，$a^2+b^2\ne0$。
\SubQ{1}{$\displaystyle\int \frac{\sin^2x\cos x}{\sin x+\cos x}\,dx$。}
\SubQ{2}{$\displaystyle\int \frac{\sin x\cos x}{1+\sin^4x}\,dx$。}
\SubQ{3}{$\displaystyle\int \frac{\sin^2x}{1+\sin^2x}\,dx$。}
\end{minipage}\par\vfill
\noindent\begin{minipage}{\textwidth}\ProblemHeading{习题}{4.52}
求下列三角函数有理式的不定积分，其中 $a,b$ 为常数，$a^2+b^2\ne0$。
\SubQ{1}{$\displaystyle\int \frac{\sin x\cos x}{\sin x+\cos x}\,dx$。}
\SubQ{2}{$\displaystyle\int \frac1{\sin^4x\cos^4x}\,dx$。}
\SubQ{3}{$\displaystyle\int \frac1{2\sin x+\sin2x}\,dx$。}
\end{minipage}\par\vfill
\newpage
\noindent\begin{minipage}{\textwidth}\ProblemHeading{习题}{4.53}
求下列三角函数有理式的不定积分，其中 $a,b$ 为常数，$a^2+b^2\ne0$。
\SubQ{1}{$\displaystyle\int \frac{\cos x}{a\sin x+b\cos x}\,dx$。}
\end{minipage}\par\vfill
